% TOP_PROBLEM: {"id": 3311, "title": "Domination in plane triangulations", "category_id": 3, "category": {"id": 3, "name": "graph_theory", "display_name": "Graph Theory", "description": "Problems involving graphs, networks, and their properties.", "slug": "graph-theory", "order_index": 3, "created_at": "2026-07-31T15:26:25.671Z"}, "status": "open", "problem_number": "OPG-36922", "set_id": 12}
% FIRST_POSTED: 2026-10-02
% ATTEMPT_STATUS: unresolved
% UnsolvedMath ID 3311; source catalogue code OPG-36922
% !TeX program = lualatex
% GPT-6 Astra Ultra research attempt, 2026-10-02; independently checked partial work.
% Completion estimate: no numerical percentage assigned; see final status and literature audit.
\documentclass[11pt,a4paper]{article}
\usepackage[margin=25mm]{geometry}
\usepackage{fontspec}
\setmainfont{lmroman10-regular.otf}[BoldFont=lmroman10-bold.otf,
 ItalicFont=lmroman10-italic.otf,BoldItalicFont=lmroman10-bolditalic.otf]
\usepackage{amsmath,amssymb,mathtools}
\usepackage{unicode-math}
\setmathfont{latinmodern-math.otf}
\AtBeginDocument{\let\setminus\smallsetminus}
\usepackage{xurl,hyperref}
\hypersetup{colorlinks=true,urlcolor=blue}
\setcounter{secnumdepth}{0}
\setlength{\parindent}{0pt}
\setlength{\parskip}{5pt}
\setlength{\emergencystretch}{3em}
\begin{document}
\section*{Domination in plane triangulations}\label{3311}
{\small UnsolvedMath ID 3311 · OPG-36922 · Graph Theory · OpenGarden}
\subsection{Definitions and mathematical statement}
A plane triangulation is a finite connected simple graph embedded in
the plane without crossings such that the boundary of every face,
including the unbounded face, is a triangle. The embedding is part
of the description. Let $G=(V,E)$ be such a graph and put $n=|V|$.
For $v\in V$, its closed neighbourhood is
$N[v]=\{v\}\cup\{u:uv\in E\}$.

A dominating set is a subset $S\subseteq V$ satisfying
$\bigcup_{v\in S}N[v]=V$. Thus every vertex outside $S$ has a neighbour
in $S$, while vertices inside $S$ already dominate themselves.
The domination number $\gamma(G)$ is the minimum cardinality of a
dominating set; no independence or connectedness condition on $S$
is imposed.

The conjecture asks whether there exists an absolute integer $N_0$
such that every plane triangulation with $n\ge N_0$ has
\[
                             \gamma(G)\le\frac n4,
\]
or equivalently $\gamma(G)\le\lfloor n/4\rfloor$. The threshold $N_0$
may be unspecified, but must be independent of the triangulation.
The quantifier allows finitely many small exceptions and requires an
exact quarter bound at every sufficiently large order.

Replacing the target by $(1/4+o(1))n$ would be an asymptotic weakening,
since a small positive error can exceed one vertex. Likewise,
allowing additional vertices on edges or inside faces changes the
instance and is not permitted. The task is to dominate the original
triangulation by selecting its original vertices.
\subsection{Short English statement}
Does every sufficiently large plane triangulation have a dominating set containing at most one quarter of its vertices?
\subsection{Research attempt}
\paragraph{Scope and process.}
GPT-6 Astra Ultra, 2 October 2026. The canonical formulation above is retained. This attempt uses a primary-source literature audit, independent restricted proofs, and adversarial checks. Reproved results are not claimed as new. Numerical tests below are reproducible in \texttt{scripts/check\_attempt\_batch03\_extremal\_2026\_10\_02.py}.

\paragraph{Literature and goal of the construction.}
Matheson and Tarjan's question and their examples are recorded in [S2]. The examples begin with disjoint copies of $K_4$ whose inner vertices remain of degree three after triangulating. We make that construction completely explicit, obtaining exact equality $\gamma=n/4$ for arbitrarily large triangulations, and also check an infinite family where the conjectured upper bound holds easily. This verifies sharpness and boundary behaviour; it does not provide the universal upper bound. The general bound $2n/7$ for $n>10$ proved by Christiansen, Rotenberg and Rutschmann [S3] is stronger than the older $n/3$ bound, but remains above the requested $n/4$.

\paragraph{An explicit triangulated insertion.}
Start with a plane $K_4$: a boundary triangle and one vertex inside it, designated private. Suppose a triangular face $abc$ contains no private vertex on its boundary. Inside that face draw a smaller triangle $ABC$, put a new private vertex $D$ inside it, and join $D$ to $A,B,C$. Add six edges across the annulus:
\[
 aA,\ bB,\ cC,\ aB,\ bC,\ cA.
\]
The old face is replaced by the six annular faces
\[
 abB,\quad bcC,\quad caA,\quad aAB,\quad bBC,\quad cCA,
\]
and the three inner faces $ABD,BCD,CAD$. This list explicitly provides a noncrossing triangulation of the region. All vertices are distinct except for the indicated shared face boundaries; no loops or parallel edges appear. The operation adds four vertices and twelve edges, maintaining $|E|=3|V|-6$.

The new annular faces have no private boundary vertex, so one is available for the next insertion. No future insertion need touch any private vertex. After $t$ copies, the graph has $4t$ vertices split into $t$ disjoint four-vertex blocks, each spanning its original $K_4$. The private vertex $z_i$ of block $i$ has closed neighbourhood exactly that block. Every dominating set must therefore meet each of these $t$ pairwise disjoint closed neighbourhoods, proving $\gamma\ge t$. Conversely choose one nonprivate vertex from each block. It is adjacent to the other three vertices of its $K_4$, so these $t$ choices dominate the whole graph. Hence
\[
 |V|=4t,\qquad \gamma=t.
\]
Additional edges between blocks cannot defeat the lower bound because they never meet the private vertices. This is the point that an unspecified completion to a triangulation would need to check.

\paragraph{A different family and the small exceptions.}
Take a cycle $C_m$, $m\ge4$, and two nonadjacent vertices $u,v$, each joined to every cycle vertex. Embed the cycle on an equator, one pole in each hemisphere. Its $2m$ faces are exactly $ux_ix_{i+1}$ and $vx_ix_{i+1}$, including whichever face is selected as the outer face in a plane drawing. Thus this bipyramid is a simple plane triangulation with $n=m+2$ vertices.

The two poles dominate all vertices, so $\gamma\le2$. No vertex is universal: each pole misses the other, while a cycle vertex has degree four and $n\ge6$. Therefore $\gamma=2$. The quarter bound holds for $n\ge8$ and fails for $n=6,7$. For $n=6$ this is the octahedron. These finite exceptions explain why the catalogue quantifies over sufficiently large order; they do not contradict that formulation.

\paragraph{Verification and why the extremal example does not prove the upper bound.}
The script carries out ten successive $K_4$ insertions. It checks the face lists, that each edge occurs on two faces, Euler's identity, all private closed neighbourhoods, and the exhibited dominating set. For the first three block counts it also finds the domination number by exhaustive search. It separately checks bipyramids of small orders. The construction deliberately preserves disjoint private neighbourhoods, which forces one choice per four vertices. An arbitrary triangulation need not split into such blocks, and greedily choosing a high-degree vertex can leave an induced remainder that is no longer a triangulation. Consequently these examples do not justify an induction using the same conjectured bound on the remainder.

\paragraph{Final status.}
Exact equality occurs in a fully specified infinite family, and all sufficiently large bipyramids satisfy the bound. The first missing statement is a domination construction of size at most $\lfloor n/4\rfloor$ for every sufficiently large triangulation. The universal target is \textbf{UNRESOLVED} by this attempt.

\subsection{Sources}
\begin{itemize}
\item[S1] ULAM.AI and the UnsolvedMath Contributors, supplied problems.json, record 3311 (OPG-36922), OpenGarden collection. Catalogue identity and source transcription; CC BY 4.0.
\url{https://huggingface.co/datasets/ulamai/UnsolvedMath}.
\item[S2] Open Problem Garden, Domination in plane triangulations. Original problem and discussion; the mathematical formulation below is expanded from the supplied source record.
\url{http://www.openproblemgarden.org/op/domination_in_plane_triangulations}.

\item[S3] A. B. G. Christiansen, E. Rotenberg and D. Rutschmann, Triangulations Admit Dominating Sets of Size $2n/7$, SODA 2024, 1194--1240; arXiv:2310.11254.
\url{https://doi.org/10.1137/1.9781611977912.47}.
\end{itemize}
\end{document}
